%%%%%%%% ICML 2026 EXAMPLE LATEX SUBMISSION FILE %%%%%%%%%%%%%%%%%

\documentclass{article}

% Recommended, but optional, packages for figures and better typesetting:
\usepackage{microtype}
\usepackage{graphicx}
\usepackage{subcaption}
\usepackage{booktabs} % for professional tables

% hyperref makes hyperlinks in the resulting PDF.
% If your build breaks (sometimes temporarily if a hyperlink spans a page)
% please comment out the following usepackage line and replace
% \usepackage{icml2026} with \usepackage[nohyperref]{icml2026} above.
\usepackage{hyperref}


% Attempt to make hyperref and algorithmic work together better:
\newcommand{\theHalgorithm}{\arabic{algorithm}}

% Use the following line for the initial blind version submitted for review:
\usepackage{icml2026}

% For preprint, use
% \usepackage[preprint]{icml2026}

% If accepted, instead use the following line for the camera-ready submission:
% \usepackage[accepted]{icml2026}

\usepackage{amsmath}
\usepackage{amssymb}
\usepackage{mathtools}
\usepackage{amsthm}


% if you use cleveref..
\usepackage[capitalize,noabbrev]{cleveref}

%%%%%%%%%%%%%%%%%%%%%%%%%%%%%%%%
% THEOREMS
%%%%%%%%%%%%%%%%%%%%%%%%%%%%%%%%
\theoremstyle{plain}
\newtheorem{theorem}{Theorem}[section]
\newtheorem{proposition}[theorem]{Proposition}
\newtheorem{lemma}[theorem]{Lemma}
\newtheorem{corollary}[theorem]{Corollary}
\theoremstyle{definition}
\newtheorem{definition}[theorem]{Definition}
\newtheorem{assumption}[theorem]{Assumption}
\theoremstyle{remark}
\newtheorem{remark}[theorem]{Remark}

% Todonotes is useful during development; simply uncomment the next line
%    and comment out the line below the next line to turn off comments
%\usepackage[disable,textsize=tiny]{todonotes}
\usepackage[textsize=tiny]{todonotes}

% The \icmltitle you define below is probably too long as a header.
% Therefore, a short form for the running title is supplied here:
\icmltitlerunning{Enhanced Pretraining for Molecular GNN Transformers}

\begin{document}

\twocolumn[
  \icmltitle{GeAT: Enhanced Multi-task Pretraining for \\
    Molecular GNN Transformers (ICML 2026)}

  % It is OKAY to include author information, even for blind submissions: the
  % style file will automatically remove it for you unless you've provided
  % the [accepted] option to the icml2026 package.

  % List of affiliations: The first argument should be a (short) identifier you
  % will use later to specify author affiliations Academic affiliations
  % should list Department, University, City, Region, Country Industry
  % affiliations should list Company, City, Region, Country

  % You can specify symbols, otherwise they are numbered in order. Ideally, you
  % should not use this facility. Affiliations will be numbered in order of
  % appearance and this is the preferred way.
  \icmlsetsymbol{equal}{*}

  \begin{icmlauthorlist}
    \icmlauthor{Mu He}{equal,yyy}
    \icmlauthor{Jinpeng Wang}{equal,yyy}
  \end{icmlauthorlist}

  \icmlaffiliation{yyy}{School of Software, Tsinghua University, Beijing, China}

  \icmlcorrespondingauthor{Mu He}{plgydamumu@fomail.com}
  \icmlcorrespondingauthor{Feihong Zhang}{zhangfeihong20@126.com}
  \icmlcorrespondingauthor{Jinpeng Wang}{TODO}

  % You may provide any keywords that you find helpful for describing your
  % paper; these are used to populate the "keywords" metadata in the PDF but
  % will not be shown in the document
  \icmlkeywords{Machine Learning, Deep Learning, GNN, Molecular Prediction, ICML}

  \vskip 0.3in
]

% this must go after the closing bracket ] following \twocolumn[ ...

% This command actually creates the footnote in the first column listing the
% affiliations and the copyright notice. The command takes one argument, which
% is text to display at the start of the footnote. The \icmlEqualContribution
% command is standard text for equal contribution. Remove it (just {}) if you
% do not need this facility.

% Use ONE of the following lines. DO NOT remove the command.
% If you have no special notice, KEEP empty braces:
\printAffiliationsAndNotice{}  % no special notice (required even if empty)
% Or, if applicable, use the standard equal contribution text:
% \printAffiliationsAndNotice{\icmlEqualContribution}

\begin{abstract}
  This paper introduces GeAT (Graph edgewise-Attention Transformer), a Transformer-like molecular graph model architecture for molecular property prediction, along with a multi-task pre-training framework for molecular graph models. GeAT's pre-training incorporates multiple self-supervised learning tasks with a two-stage multi-task weighting strategy. This pre-training framework enables the model to learn general features from large-scale unlabeled molecular datasets. GeAT models edges as aggregation functions and computes attention scores for different bond types through an explicit edge-aware attention mechanism, allowing the model to learn bond-type-specific interaction patterns. GeAT achieves state-of-the-art performance across multiple benchmark datasets, demonstrating strong capabilities in molecular representation learning.
\end{abstract}

\section{Introduction}

In recent years, graph neural networks (GNNs) have shown strong potential in molecular representation, since molecules can be naturally modeled as graphs. 

Pioneered by Hu et al., a series of works have leveraged large-scale unlabeled molecular datasets to design pre-training tasks. These tasks enable GNNs to learn general molecular representations from unlabeled data. When fine-tuned on downstream tasks such as property prediction, pre-trained GNNs achieve better performance.

This paper enhances existing pre-training frameworks for molecular graph models.
It introduces a series of progressively designed pre-training tasks:
\begin{itemize}
\item Node attribute prediction.
\item Masked node type prediction.
\item Functional group recognition.
\item Molecular scaffold contrastive learning.
\item Graph-level contrastive learning.
\item Augmented contrastive learning.
\end{itemize}

We employ a two-stage multi-task learning weighting strategy for these tasks.  

In the first stage, we use Uncertainty Weighting to align the scales of different tasks and estimate their relative difficulties. The weights learned in this stage also serve as initialization for the second stage.

In the second stage, we apply GradNorm. GradNorm normalizes task gradients to balance their contributions during training. This enables stable and balanced convergence of all task losses.

These tasks help the model learn richer molecular representations.
As a result, we witness performance improvements on downstream tasks. 

In addition, we propose GeAT, a Transformer-like architecture tailored for molecular tasks. 

Unlike prior work by Hu et al., which represents edges in molecular graphs as embedding vectors, GeAT treats edges as aggregation functions. It introduces an edge-aware attention mechanism that performs sparse routing based on edge information, allowing different attention blocks to compute interaction weights between neighboring nodes. 

The model combines edge-attention layers and global-attention layers to capture interaction patterns at multiple scales. A small Mixture-of-Experts (MoE) module is used as the final output layer to produce molecular representations. When paired with the pre-training framework described earlier, GeAT achieves state-of-the-art results on downstream classification and regression tasks in the public MoleculeNet benchmark.

For tasks involving 3D molecular information, we also designed additional self-supervised pre-training tasks to help the model learn geometric-aware representations. These include regression tasks for predicting the bond angles of neighboring node triplets and the dihedral angles of quadruplets in the molecular graph.  

We evaluated the effectiveness of these 3D-aware tasks using multiple 3D-related molecular properties from the open-source NablaDFT dataset as downstream prediction targets. The results confirm the benefits of incorporating 3D geometric information during pre-training.

In summary, our key contributions are as follows:
\begin{itemize}
\item We propose GeAT, a novel variant of GAT that models edges as aggregation functions over node embeddings and employs attention layers at multiple scales to enable richer semantic modeling.
\item We enhance the molecular GNN pre-training framework with a series of progressively structured self-supervised tasks—covering both semantic and 3D-aware objectives—and optimize them using a two-stage multi-task weighting strategy: Uncertainty Weighting for initial task balancing, followed by GradNorm to stabilize convergence. This enables effective learning of general molecular representations from large-scale unlabeled data.
\item Based on GeAT and the enhanced pre-training strategy, we provide a reproducible training and fine-tuning pipeline, along with readily usable pre-trained model weights. This setup achieves state-of-the-art performance on multiple downstream benchmarks, including MoleculeNet.
\end{itemize}

\section{Related Work}
\subsection{GNN architectures}
\subsubsection{Traditional GNNs}
\subsubsection{Molecular GNNs}
\subsection{Molecular GNN Pre-training}

\section{Preliminary}
\subsection{Molecular GNN}
\subsection{Multi-task Learning Weight Strategies}
\subsubsection{Uncertainty Weighting}
\subsubsection{GradNorm}

\section{GeAT Pre-training Framework}
This section presents the GeAT model architecture and the pre-training framework, which includes a series of self-supervised tasks—both semantic and 3D-aware—as well as fine-tuning configurations for downstream tasks.

\subsection{Edge-wise Attention}
The edge-wise attention mechanism in GeAT is based on a key principle: different chemical bond types—such as single, double, and aromatic bonds—require distinct interaction patterns. 

Unlike conventional molecular GNNs that encode edges as embedding vectors for computation, GeAT encodes each bond type as an attention-based aggregation function. 

Formally, let $\operatorname{T}(i,j)$ denote the type attribute of the undirected edge $(i,j)$, the $k$-th layer of edge-wise attention blocks for node $i$ is: 
\begin{equation}
e_{ij}^{(k)} = \operatorname{LeakyReLU}\left(\operatorname{Aggr}_{\operatorname{T}(i,j)}\left(\mathbf{W}_{Q}\mathbf{h}_i^{(k)}, \mathbf{W}_{K}\mathbf{h}_j^{(k)} \right)\right),  
\end{equation}
\begin{equation}
\alpha_{ij}^{(k)} = \frac{\exp(e_{ij}^{(k)})}{\sum_{k \in \mathcal{N}(i)} \exp(e_{ik}^{(k)})},
\end{equation}
\begin{equation}
\mathbf{h}_i^{(k+1)} = \operatorname{Proj}\left( \sum_{j \in \mathcal{N}(i)} \alpha_{ij}^{(k)} \, \mathbf{W}_V \mathbf{h}_j^{(k)} \right),  
\end{equation}

The $\operatorname{Aggr}$ function is typically a simple bilinear operation:
\begin{equation}
\operatorname{Aggr}_{\operatorname{T}(i,j)}\left(\mathbf{x},\mathbf{y}\right)=\mathbf{x}^\top\mathbf{A}_{\operatorname{T}(i,j)}\mathbf{y}, 
\end{equation}

or a two-layer MLP can also be used instead to increase capacity:
\begin{equation}
\begin{aligned}
\operatorname{Aggr}_{\operatorname{T}(i,j)}\left(\mathbf{x}, \mathbf{y}\right)
&= \operatorname{MLP}_{\operatorname{T}(i,j)}\!\left( [\mathbf{x} \| \mathbf{y}] \right) \\
&= \operatorname{ReLU}\left( \mathbf{w}^{(2)}_{\operatorname{T}(i,j)} \mathbf{z}_1 + b^{(2)}_{\operatorname{T}(i,j)} \right), \\
\text{where } \mathbf{z}_1 &= \operatorname{ReLU}\left( \mathbf{W}^{(1)}_{\operatorname{T}(i,j)} [\mathbf{x} \| \mathbf{y}] + \mathbf{b}^{(1)}_{\operatorname{T}(i,j)} \right),
\end{aligned}
\end{equation}

where  $[\mathbf{x}|| \mathbf{y}]$  denotes the concatenation of $\mathbf{x}$ and $\mathbf{y}$, and $\operatorname{MLP}_t$ is a dedicated multi-layer perceptron for edge type $t$.

The $\operatorname{Proj}$ module applies a projection matrix to map the output dimension to the input dimension of the next layer—especially important in our multi-head attention. Assume there are $\mathbf{H}$ heads:
\begin{equation}
e_{ij}^{(k,h)} = \operatorname{LeakyReLU}\left( \operatorname{Aggr}_{\operatorname{T}(i,j)}\left( \mathbf{W}_Q^{(h)} \mathbf{h}_i^{(k)},\, \mathbf{W}_K^{(h)} \mathbf{h}_j^{(k)} \right) \right),
\end{equation}

\begin{equation}
\alpha_{ij}^{(k,h)} = \frac{\exp\!\left(e_{ij}^{(k,h)}\right)}{\sum_{l \in \mathcal{N}(i)} \exp\!\left(e_{il}^{(k,h)}\right)}.
\end{equation}

\begin{equation}
\mathbf{h}_i^{(k+1)} = \operatorname{Proj}\left( \big\|_{h=1}^{\mathbf{H}} \sum_{j \in \mathcal{N}(i)} \alpha_{ij}^{(k,h)} \, \mathbf{W}_V^{(h)} \mathbf{h}_j^{(k)} \right),
\end{equation}

Our multi-head attention follows the GAT design: inputs are first projected into a higher-dimensional space using  $\mathbf{W}_Q,\mathbf{W}_K,\mathbf{W}_V \in \mathbf{R}^{d\times dH}$ , split across multiple attention heads, processed in parallel, and then projected back to the original dimension via $\operatorname{Proj}$. 

\subsection{GNN architectures}

\subsection{Pre-training Framework}
\subsubsection{Weighting Strategy}

\subsubsection{Pre-training Tasks}
Our pre-training framework comprises six tasks, organized into three progressive levels: node-level semantics, local substructure semantics, and graph-level semantics. For downstream tasks requiring 3D information, we further include two additional 3D-aware tasks to capture spatial semantics.

\textbf{Node attribute prediction} 

ChemProp provides a set of effective handcrafted atomic features, such as formal charge, hybridization, chirality, number of hydrogens and atomic mass. This task reconstructs these features from atomic embeddings to preserve essential node-level semantics. 

This task includes both classification and regression components.  
We use the sum of binary cross-entropy (BCE) loss for classification and mean squared error (MSE) loss for regression as the total loss:
\begin{align*}
\mathcal{L}_{\text{nodeattr}} 
&= -\frac{1}{N_{\text{cl}}} \sum_{i=1}^{N_{\text{cl}}} \left[ y_{\text{cl},i} \log(\hat{y}_{\text{cl},i}) + (1 - y_{\text{cl},i}) \log(1 - \hat{y}_{\text{cl},i}) \right] \\
&\quad + \frac{1}{N_{\text{reg}}} \sum_{j=1}^{N_{\text{reg}}} \left( y_{\text{reg},j} - \hat{y}_{\text{reg},j} \right)^2
\end{align*}

\textbf{Masked node type prediction}

\textbf{Functional group recognition}

\textbf{Molecular scaffold contrastive learning}

\textbf{Graph-level contrastive learning}

\textbf{Augmented Contrastive learning}

\subsection{Fine-tuning for Downstream Tasks}

\section{Experiments}
\subsection{Setup}
\subsubsection{Datasets}
\subsubsection{Training Configurations}
\subsection{Evaluation of Model Architecture}
\begin{table*}[t]
  \caption{Performance comparison on classification and regression tasks.
  For classification tasks, results are reported as AUC; for regression tasks, as MSE loss. All models are evaluated over 10 random seeds (0–9). The numbers in brackets are the standard deviation.}
  \begin{center}
    \begin{small}
      \begin{sc}
        \begin{tabular}{l|cccccc|ccccc}
          \toprule
          & \multicolumn{6}{c|}{(classification)} & \multicolumn{5}{c}{(regression)} \\
          & BBBP & SIDER & Clintox & BACE & Tox21 & ToxCast & FreeSolv & ESOL & Lipo & QM7 & QM8 \\
          \midrule
          Breast       & 95.9(0.2 & $\surd$ & & & & & & & & & \\
          Cleveland    & 83.3(0.6 & $\times$ & & & & & & & & & \\
          Glass2       & 61.9(1.4 & $\surd$ & & & & & & & & & \\
          Credit       & 74.8(0.5 & 78.3(0.6 & & & & & & & & & \\
          Horse        & 73.3(0.9 & 69.7(1.0 & $\times$ & & & & & & & & \\
          Meta         & 67.1(0.6 & 76.5(0.5 & $\surd$ & & & & & & & & \\
          \midrule
          GeAT(no pre) & 75.1(0.6 & 73.9(0.5 & & & & & & & & & \\
          GeAT         & 44.9(0.6 & 61.5(0.4 & $\surd$ & & & & & & & & \\
          \bottomrule
        \end{tabular}
      \end{sc}
    \end{small}
  \end{center}
  \vskip -0.1in
\end{table*}
\subsection{Evaluation of Pre-training Framework}
\subsection{Evaluation of 3D-aware Tasks}

\section{Conclusions}

\section{Limitations And Future Work}

\bibliography{example_paper}
\bibliographystyle{icml2026}
%%%%%%%%%%%%%%%%%%%%%%%%%%%%%%%%%%%%%%%%%%%%%%%%%%%%%%%%%%%%%%%%%%%%%%%%%%%%%%%
%%%%%%%%%%%%%%%%%%%%%%%%%%%%%%%%%%%%%%%%%%%%%%%%%%%%%%%%%%%%%%%%%%%%%%%%%%%%%%%
\end{document}

% This document was modified from the file originally made available by
% Pat Langley and Andrea Danyluk for ICML-2K. This version was created
% by Iain Murray in 2018, and modified by Alexandre Bouchard in
% 2019 and 2021 and by Csaba Szepesvari, Gang Niu and Sivan Sabato in 2022.
% Modified again in 2023 and 2024 by Sivan Sabato and Jonathan Scarlett.
% Previous contributors include Dan Roy, Lise Getoor and Tobias
% Scheffer, which was slightly modified from the 2010 version by
% Thorsten Joachims & Johannes Fuernkranz, slightly modified from the
% 2009 version by Kiri Wagstaff and Sam Roweis's 2008 version, which is
% slightly modified from Prasad Tadepalli's 2007 version which is a
% lightly changed version of the previous year's version by Andrew
% Moore, which was in turn edited from those of Kristian Kersting and
% Codrina Lauth. Alex Smola contributed to the algorithmic style files.
